\section{The framework}
\label{sec:framework}

Figure~\ref{fig:f1}a summarizes the offline/online pipeline; each component below is introduced
by the production failure mode it removes. Throughout, \(y_t\) is the (delayed,
infrequent) quality measurement, \(\hat{y}_t\) the soft-sensor estimate, and \(\theta\) the
label delay in samples, treated throughout as a first-class protocol parameter rather than an afterthought.

\subsection{Physics-derived features and transport-lag diagnosis}

\emph{Failure mode: complexity set by validation.} Rather than offering the model every
input at every lag and letting validation choose (the \(k = 126\) kitchen sink of
Section 5.1), the feature set is fixed by the process physics. For the debutanizer
column, three quantities follow from the vapor--liquid equilibrium of the
butane--stabilized-gasoline separation: a bubble-point estimate of the light-key
fraction obtained by inverting the bubble-point condition at the measured tray
temperature and pressure (Antoine/DIPPR-101 vapor pressures; n-hexane as the
heavy-component proxy) \citep{smith2018}; the relative volatility \(\alpha(T)\)
evaluated at the same state; and a stripping factor coupling \(\alpha(T)\) with the
reflux signal. The raw tray temperature at the transport lag is retained as a robust
backbone, giving \(k = 4\) features.

The transport lag itself is diagnosed, never assumed. A lag scan of the
target--temperature cross-correlation exposes the column's hydraulic transport delay as
a stable interior maximum (\(u_5\) at lag 15, \(r^2 \approx 0.51\) on the debutanizer); at
lag 0 the same physics appears to fail (\(R^2 \approx 0.02\)). The diagnostic posture
matters: a lag-zero failure indicts the data alignment, not the thermodynamics, and
resolving it costs one scan rather than a model-class escalation.

\subsection{Selection: blocked time-series cross-validation with the one-SE rule}

\emph{Failure mode: the selection lottery under autocorrelation and regime shift.}
Validation on adjacent or randomly shuffled samples leaks autocorrelated information
and averages away regime structure; selection then varies with the seed. We use
\(K\)-fold blocked CV (\(K = 5\)): the pool is cut into contiguous time blocks, features
are built within each segment separately (so lagged features cannot straddle a
boundary), the scaler is fit on fold-train only, and a gap of the maximum lag
separates train from validation blocks. The per-fold \(R^2\) spread, including the
strongly negative folds in Table~\ref{tab:t1}, is retained as the signal rather than
suppressed as noise.

Formally, let the pool be \(\{(x_t, y_t)\}_{t=1}^{N}\) in time order and let
\(B_1, \dots, B_K\) be its partition into \(K\) contiguous blocks of (near-)equal length.
For fold \(k\), the validation set is \(B_k\) and the training set is
\(\bigcup_{j \neq k} B_j\) with the \(\ell\) samples adjacent to each boundary of \(B_k\)
removed, where \(\ell\) is the maximum lag used by any feature; feature construction is
applied to each contiguous segment independently, so no lagged feature spans a
boundary in either direction. The fold score is the validation \(R^2\) of the pipeline
(imputer where applicable, scaler, estimator) fit on the gapped training set, and we
report \(\bar{R}^2 \pm \mathrm{SE}\) over folds together with the worst fold. Two
properties matter. The gap makes the fold estimate leakage-free under serial
dependence of range \(\leq \ell\): no validation point shares lagged inputs with any
training point. Contiguity, in turn, makes each fold a \emph{regime probe}: if the pool spans
operating regimes, at least one fold must validate on a regime the model under-saw in
training, which is the very deployment event the per-fold spread is meant to
price.

Model choice on any complexity path (feature sets; the elastic-net \(\alpha\) path of
Section 5.5) then follows the one-standard-error rule \citep{hastie2009}: take the
complexity with the best mean CV score and select the \emph{simplest} complexity whose mean
lies within one standard error of it. The rule encodes that validation cannot
distinguish the robust low-complexity model from the regime-overfit high-complexity
one, and resolves the tie toward parsimony without ever touching the held-out test
set. On SECOM this is the difference between \(R^2 \approx -1.8\) and \(R^2 \approx -10^4\) (Figure~\ref{fig:f3}).

\subsection{Drift: open-loop bias update with a corrected-residual alarm}

\emph{Failure mode: slow drift and regime offsets under delayed, infrequent labels.}
Following the open-loop case of \citep{shardt2016}, the deployed estimate carries a
causal EWMA bias term driven only by labels that have already arrived:

\[
b_t = (1 - \lambda)\, b_{t-1} + \lambda \left( y_{t-\theta} - \hat{y}_{t-\theta} \right),
\qquad
\tilde{y}_t = \hat{y}_t + b_t,
\]

with \(b\) held at \(b_0 = 0\) until the first label. \(\lambda\) is CV-selected on the pool
(test untouched; \(\lambda = 0.1\) on the debutanizer), and \(\theta\) is the documented
analyzer delay. The update is open-loop in the sense that it never feeds a controller, so the
stability concerns of closed-loop bias correction do not arise; what it buys is the
9$\times$ cross-fold SE reduction and the rescue of regime-shifted folds reported in
Section 5.2. Because a tracking update can absorb genuine drift, drift \emph{alarming} is
performed on the corrected residual \(y_t - \tilde{y}_t\) (ADWIN \citep{bifet2007adwin}): the detector fires only when the bias update can no longer keep up, which is
the event that warrants intervention.

\subsection{Calibrated uncertainty: adaptive conformal intervals under label delay}

\emph{Failure mode: uncalibrated confidence under drift.} Split conformal prediction
\citep{papadopoulos2002} turns held-out absolute residuals
\(s_i = |y_i - \tilde{y}_i|\) into a distribution-free interval
\(\tilde{y}_t \pm Q_{1-\alpha}(s)\), but its guarantee requires exchangeability,
which is exactly what drift destroys (Section 5.3). We therefore apply adaptive conformal
inference \citep{gibbs2021aci}: after each label, the working miscoverage level is
updated by

\[
\alpha_{t+1} = \alpha_t + \gamma\,(\alpha - \mathrm{err}_t),
\qquad
\mathrm{err}_t = \mathbf{1}\{ y_t \notin C_t \},
\]

with target \(\alpha = 0.10\), step \(\gamma = 0.05\), a rolling score window of 200, and
\(\alpha_t\) intentionally left unclipped; the score quantile handles the
extremes, matching the original construction. The interval is formed around the
\emph{bias-corrected} estimate, so the conformal layer prices residual uncertainty after
the drift mechanism of \S 3.3 has done its work; the composition, not either piece
alone, produces the regime-uniform coverage of Table~\ref{tab:t2}. EnbPI \citep{xu2021enbpi} is
carried as a comparator and behaves as its FIFO residual ensemble predicts.

Label delay interacts with conformal bookkeeping in a way that is easy to get wrong:
when a label for sample \(i\) arrives \(\theta\) samples late, its coverage indicator must
be evaluated against the interval \emph{issued at prediction time} \(C_i\), not the interval
the (since-adapted) state would issue now. The online implementation (\S 3.6) enforces this
structurally by recording \((\hat{y}_i, \tilde{y}_i, C_i)\) per sample identifier and
scoring arrivals against the stored tuple.

\subsection{Transfer: functional scale-bias migration, composed with the online layer}

\emph{Failure mode: data scarcity on a new operating regime.} When a sibling regime of the
same process exists, the source sensor is migrated rather than retrained. Among the
scale-bias correction family, the functional form of \citep{yan2011},

\[
\tilde{y}^{\mathrm{new}}(x) = s(x)\, \hat{y}^{\mathrm{src}}(x) + \delta(x),
\qquad \delta \sim \mathcal{GP}(0, k),
\]

keeps the source model's structure while a zero-mean GP bias repairs the relationship
change locally, and it supplies calibrated posterior intervals at no extra cost. Two
scoping notes carry into the results. First, parameter-level migration requires a
shared input space; across genuinely different processes (debutanizer $\to$ TEP) it is
mathematically inapplicable, and what transfers is the \emph{methodology}: the full
\S 3.1--3.4 recipe, applied verbatim. Second, for a \emph{linear} source the per-input affine
re-scaling of \citep{luo2015} absorbs the source weights entirely and provably
collapses to from-scratch regression; Section 5.4 shows the empirical face of that
identity. Migration is offline; deployed migrated sensors still run the \S 3.3 online
update, and this two-layer composition is load-bearing (the bare-migration control of
Section 5.4 caps below threshold).

\subsection{Online implementation: estimation and delayed-label reconciliation}

\emph{Failure mode: a sensor whose calibration depends on delayed laboratory
results cannot be implemented as a fixed function.} The implemented object is a frozen
point model plus a small calibration state $(b_t, \alpha_t, \text{score window})$
that evolves as laboratory results arrive. The implementation (Figure~\ref{fig:f1}b)
therefore separates two activities running at different rates: a fast
\emph{estimation} path, executed at the process scan rate, which reads the
calibration state and issues $\tilde{y}_t$ with its interval $C_t$; and a slow
\emph{reconciliation} path, executed whenever a laboratory result arrives, which
updates $b_t$, $\alpha_t$, and the residual window, pairing each result with the
estimate and interval recorded when its sample was taken (\S 3.4). The calibration
state is persisted so that the sensor survives a restart without losing its
correction, while the point model itself is never modified online. The computational
cost is negligible by design (a linear map, an exponentially weighted average,
and a quantile lookup, each $O(1)$ per laboratory result), so the measured
99th-percentile latency of 1.97~ms per estimate (Section~5.6) leaves two orders of
magnitude of margin inside a typical 200~ms analyzer-cycle budget. Fast, auditable
computation is a design consequence of the linear scope, not an optimization
afterthought.

% --- figures first referenced in this section ---

\begin{figure*}[t]
\centering
\includegraphics[width=\textwidth]{F1_architecture}
\caption{\textbf{Framework and online implementation.} (a) The offline/online pipeline: physics-derived features (bubble-point estimate, relative volatility $\alpha(T)$, stripping factor) feed blocked time-series CV with one-SE selection; the frozen linear sensor is corrected online by an open-loop EWMA bias update at the documented analyzer delay $\theta$, and adaptive conformal intervals are computed on the corrected residuals; drift alarming acts on the corrected residual, and migration (offline) initializes sensors for sibling regimes. (b) The online implementation as a fast estimation path and a slow reconciliation path over a small calibration state ($b_t$, $\alpha_t$, residual window): each estimate is issued with its interval and recorded; each delayed laboratory result updates the calibration state and is scored against the interval recorded when its sample was estimated (dashed), not the since-adapted current one.}
\label{fig:f1}
\end{figure*}
